\begin{frame}{실습: 가상 발표를 \emph{리뷰어의 눈}으로 재실행}
  \small
  \begin{enumerate}
    \item \textbf{발표팀 재현} — Pendulum-v1, PPO(Ch11.2 코드), 시드 42,
      $\gamma=0.99$, $\lambda=0.95$, 300$\times$400(12만 스텝) $\to$
      처음 10: $-805.3$, 마지막 10: $-742.0$. 발표팀의 숫자 재현.
    \item \textbf{패턴 1 검증 — 시드 2개 추가} — 시드 123, 7로 재실행
      $\to$ 시드별 마지막 10 이동평균 표. 편차를 보고 ``알고리즘의
      차이''인가 ``시드의 우연''인가를 판정.
    \item \textbf{패턴 2 검증 — 학습 스텝 확장} — 12만 $\to$ 24만 스텝,
      이동평균이 포화되는지 $\to$ ``절반도 안 왔다''(11.2)를 숫자로
      확인.
    \item \textbf{패턴 3 검증 — 보상 스케일링 민감성} — 1/8 스케일링 +
      $\gamma=0.95$로 재실행 $\to$ 스케일링이 \emph{결론의 방향}을
      바꾸는지.
    \item \textbf{패턴 5 검증 — 베이스라인 실측} — 무작위 정책
      300 에피소드 $\to$ 평균 약 $-1239$.
    \item \textbf{리뷰 작성} — 실습에서 나온 숫자를 근거로 3점 리뷰를
      실제 템플릿(나중)으로 완성.
  \end{enumerate}
  \vskip0.4em
  {\footnotesize 이 6단계 = ``리뷰어는 발표팀의 실험을 \textbf{자기
  손으로} 재실행할 수 있는 사람''이라는 이 절의 핵심 태도를 코드
  레벨에서 보여줌. ``이 학습 곡선, 시드를 바꿔도 같은 모양이 나올까요?''
  라는 질문은 \textbf{리뷰어가 실제로 다시 돌려본 적}이 있을 때 가장
  날카롭다.}
\end{frame}
